\documentclass[11pt,a4paper]{article}
\usepackage[utf8]{inputenc}
\usepackage{booktabs}
\usepackage{longtable}
\usepackage{geometry}
\usepackage{hyperref}
\usepackage{amsmath}
\usepackage{amssymb}
\usepackage{graphicx}
\usepackage{float}
\geometry{margin=1in}

\begin{document}

\begin{center}
\Large\textbf{Supplementary Materials}
\end{center}

\tableofcontents
\newpage

% ============================================================================
% S0: TARGET SEVERITY DISTRIBUTION DERIVATION
% ============================================================================
\section{S0: Target Severity Distribution Derivation}

\subsection{S0.1 Overview}

We derived target severity distributions from DDInter, a literature-curated DDI database with peer-reviewed severity annotations (Table~\ref{tab:derivation}).

\begin{table}[h]
\centering
\scriptsize
\caption{DDInter Evidence for Target Severity Distribution}
\label{tab:derivation}
\begin{tabular}{@{}p{3.0cm}p{4.5cm}p{5.0cm}@{}}
\toprule
\textbf{Source} & \textbf{Data} & \textbf{Severity Levels} \\
\midrule
\textbf{DDInter Database} & 37,164 matched CV/AT DDI pairs & Literature-validated labels: Major (22.1\%), Moderate (73.5\%), Minor (4.4\%). No Contraindicated category. \\
\bottomrule
\end{tabular}
\end{table}

\subsection{S0.2 DDInter Database}

DDInter is a publicly accessible, literature-curated DDI database with severity annotations derived from scientific literature, package inserts, and regulatory labels (Xiong et al., \textit{Nucleic Acids Research}, 2022). \textbf{Importantly, DDInter provides drug pairs and severity labels but does not include interaction description text.} We matched 37,164 DDI pairs between our DrugBank cardiovascular/antithrombotic dataset (which provides descriptions) and DDInter (which provides severity labels) using exact drug name matching. The matched subset showed: Major 22.1\%, Moderate 73.5\%, Minor 4.4\%. DDInter does not distinguish a Contraindicated category from Major.

\subsection{S0.3 Target Distribution Selection}

We aligned our targets with DDInter's matched CV/AT distribution. Since DDInter lacks a Contraindicated category, we used FAERS pharmacovigilance data to split DDInter's Major (22\%) into Contraindicated and Major (Table~\ref{tab:consensus}).

\textbf{Rationale for each target:}
\begin{itemize}
\item \textbf{Contraindicated + Major (22\%):} Matches DDInter's Major (22.1\%). To split this into Contraindicated vs.\ Major, we analyzed FAERS multi-drug adverse event reports where death/life-threatening outcomes comprised approximately 20\% of serious events. Applying this ratio: Contraindicated $\approx$ 4\%, Major $\approx$ 18\%.
\item \textbf{Moderate (74\%):} Preserved from DDInter's 73.5\%.
\item \textbf{Minor (4\%):} Preserved from DDInter's 4.4\%.
\end{itemize}

\begin{table}[h]
\centering
\small
\caption{DDInter Evidence and Selected Targets}
\label{tab:consensus}
\begin{tabular}{@{}lcc@{}}
\toprule
\textbf{Severity} & \textbf{DDInter (matched)} & \textbf{Selected Target} \\
\midrule
Contraindicated & Major: 22.1\% & \textbf{4\%} \\
Major & (combined) & \textbf{18\%} \\
Moderate & 73.5\% & \textbf{74\%} \\
Minor & 4.4\% & \textbf{4\%} \\
\bottomrule
\end{tabular}
\end{table}

\newpage

% ============================================================================
% S1: CLASSIFICATION METHODS
% ============================================================================
\section{S1: Classification Methods}

This section provides detailed equations and algorithms for the classification approaches described in the main text: (1) Zero-Shot BART-MNLI, (2) Confidence-Weighted, (3) Rule-Based Classifier, and (4) Evidence-Based Classifier (Section S2).

\subsection{S1.1 Zero-Shot BART-MNLI}

The Zero-Shot method uses \texttt{facebook/bart-large-mnli}, a 400M-parameter natural language inference model. Severity classification is framed as textual entailment: each DDI description serves as the \textit{premise}, tested against four \textit{hypotheses} (``This interaction is contraindicated/major/moderate/minor''). The model outputs entailment probabilities; the highest-probability class is assigned:
\begin{equation}
\hat{y} = \arg\max_{c \in \mathcal{C}} P(\text{entailment} \mid \text{description}, c)
\end{equation}
where $\mathcal{C} = \{\text{Contraindicated, Major, Moderate, Minor}\}$.

\subsection{S1.2 Confidence-Weighted Adjustment}

The Confidence-Weighted method extends Zero-Shot BART-MNLI by downgrading high-severity predictions when model confidence is low ($<$65\%):
\begin{equation}
\text{Severity}_{\text{adj}} = 
\begin{cases}
\text{downgrade by 1 level} & \text{if } \hat{y} \in \{\text{Contraindicated, Major}\} \land c < 0.65 \\
\hat{y} & \text{otherwise}
\end{cases}
\end{equation}
where $\hat{y}$ is the predicted severity and $c$ is the model's confidence score (maximum entailment probability).

\subsection{S1.3 Rule-Based Classifier}

The Rule-Based classifier assigns severity based on clinical keyword pattern matching in DDI descriptions. We matched 37,164 DDI pairs between DrugBank (which provides interaction descriptions) and DDInter (which provides expert-validated severity labels but no descriptions). Keyword weights $w_k$ were \textbf{empirically derived} using log-likelihood ratios.

\subsubsection{S1.3.1 Train/Test Split}

We split the 37,164 matched DDI pairs into training (70\%, n=26,014) and held-out test (30\%, n=11,150) sets using stratified sampling to preserve severity class proportions.

\subsubsection{S1.3.2 Empirical Weight Derivation}

For each clinical keyword $k$, we computed:
\begin{equation}
w_k = \ln\left(\frac{P(k|\text{Major})}{P(k|\text{Moderate})}\right)
\end{equation}
where $P(k|\text{Level})$ is the frequency of keyword $k$ appearing in DrugBank descriptions for DDI pairs labeled with that severity level in DDInter. Positive weights indicate keywords predictive of Major severity; negative weights indicate keywords predictive of Moderate or Minor severity.

\begin{table}[H]
\centering
\caption{Empirically-Derived Keyword Weights (from training set)}
\label{tab:empirical_weights}
\begin{tabular}{@{}lcccc@{}}
\toprule
\textbf{Keyword} & \textbf{Weight} & \textbf{LR} & \textbf{Major \%} & \textbf{Mod \%} \\
\midrule
\multicolumn{5}{l}{\textit{Positive weights (predictive of Major):}} \\
prolongation & +1.45 & 4.25$\times$ & 11.6\% & 2.7\% \\
bleeding & +1.03 & 2.80$\times$ & 25.9\% & 9.2\% \\
hemorrhage & +0.63 & 1.88$\times$ & 9.2\% & 5.0\% \\
anticoagulant & +0.57 & 1.76$\times$ & 3.3\% & 1.8\% \\
risk / severity & +0.33 & 1.39$\times$ & 49.7\% & 35.3\% \\
increased & +0.20 & 1.22$\times$ & 64.0\% & 52.0\% \\
\midrule
\multicolumn{5}{l}{\textit{Negative weights (predictive of Moderate/Minor):}} \\
decrease & $-$0.30 & 0.74$\times$ & 25.7\% & 35.3\% \\
serum & $-$0.31 & 0.73$\times$ & 13.3\% & 18.5\% \\
therapeutic & $-$1.20 & 0.30$\times$ & 3.3\% & 11.8\% \\
efficacy & $-$1.27 & 0.28$\times$ & 3.8\% & 14.5\% \\
antihypertensive & $-$1.49 & 0.23$\times$ & 0.5\% & 2.3\% \\
reduce & $-$1.61 & 0.20$\times$ & 0.3\% & 1.7\% \\
\bottomrule
\end{tabular}
\end{table}

LR = Likelihood Ratio = $P(k|\text{Major}) / P(k|\text{Moderate})$. Keywords with LR $>$ 1 are more frequent in Major interactions. Weight = $\ln(\text{LR})$.

\subsubsection{S1.3.3 Algorithm: Keyword Weight Derivation}

\textbf{Algorithm 1: Empirical Keyword Weight Derivation}

\textit{Input:} $\mathcal{D}_{\text{train}}$ (training set of DrugBank description, DDInter severity pairs), $\mathcal{K}$ (candidate keywords)

\textit{Output:} $\mathcal{W}$ (dictionary of keyword $\rightarrow$ weight)

\begin{enumerate}
\item Count Major and Moderate instances: $n_{\text{Major}}$, $n_{\text{Moderate}}$
\item For each keyword $k \in \mathcal{K}$:
\begin{itemize}
\item Count occurrences: $c_{\text{Major}}$, $c_{\text{Moderate}}$
\item If $c_{\text{Major}} + c_{\text{Moderate}} \geq 50$:
\begin{itemize}
\item $P(k|\text{Major}) \gets c_{\text{Major}} / n_{\text{Major}}$
\item $P(k|\text{Moderate}) \gets c_{\text{Moderate}} / n_{\text{Moderate}}$
\item $\mathcal{W}[k] \gets \ln(P(k|\text{Major}) / P(k|\text{Moderate}))$
\end{itemize}
\end{itemize}
\item Return $\mathcal{W}$
\end{enumerate}

\subsubsection{S1.3.4 Algorithm: Severity Classification}

\textbf{Algorithm 2: Rule-Based Severity Classification}

\textit{Input:} $\mathcal{W}$ (keyword weights), $\mathcal{S}_{\text{train}}$ (training scores), $d$ (DDI description)

\textit{Output:} severity $\in$ \{Contraindicated, Major, Moderate, Minor\}

\textbf{Step 1: Compute severity score}
\begin{equation}
S = \sum_{k \in \mathcal{W}} \mathcal{W}[k] \cdot \mathbb{1}[k \in d]
\end{equation}

\textbf{Step 2: Compute percentile thresholds from training scores}
\begin{itemize}
\item $P_{96} \gets \text{percentile}(\mathcal{S}_{\text{train}}, 96)$ (top 4\%)
\item $P_{78} \gets \text{percentile}(\mathcal{S}_{\text{train}}, 78)$ (top 22\%)
\item $P_{4} \gets \text{percentile}(\mathcal{S}_{\text{train}}, 4)$ (bottom 4\%)
\end{itemize}

\textbf{Step 3: Classify based on score}
\begin{equation}
\text{severity} = 
\begin{cases}
\text{Contraindicated} & \text{if } S \geq P_{96} \\
\text{Major} & \text{if } P_{78} \leq S < P_{96} \\
\text{Minor} & \text{if } S \leq P_{4} \\
\text{Moderate} & \text{otherwise}
\end{cases}
\end{equation}

\subsubsection{S1.3.5 Severity Score Calculation}

For each DDI description, the severity score is computed:
\begin{equation}
S = \sum_{k \in \text{keywords}} w_k \cdot \mathbb{1}[k \in \text{description}]
\end{equation}
where $\mathbb{1}[\cdot]$ is the indicator function.

\subsubsection{S1.3.6 Percentile-Based Thresholds}

Final classification uses percentile thresholds derived from the training set score distribution, aligned to DDInter's severity distribution:

\begin{itemize}
\item \textbf{Contraindicated:} Top 4\% of scores ($\geq$ 96th percentile)
\item \textbf{Major:} Next 18\% of scores (78th--96th percentile)
\item \textbf{Moderate:} Middle 74\% (4th--78th percentiles)
\item \textbf{Minor:} Bottom 4\% of scores ($\leq$ 4th percentile)
\end{itemize}

\begin{equation}
\text{Severity} = 
\begin{cases}
\texttt{Contraindicated} & S \geq P_{96} \\
\texttt{Major} & P_{78} \leq S < P_{96} \\
\texttt{Minor} & S \leq P_{4} \\
\texttt{Moderate} & \text{otherwise}
\end{cases}
\end{equation}

\subsubsection{S1.3.7 Validation Results}

\textbf{DDInter Validation Mapping:} For validation against DDInter (which uses only 3 levels), we map Contraindicated $\rightarrow$ Major, yielding a combined Major = 22\% matching DDInter's distribution.

\textbf{Test Set Performance (n=11,150):}
\begin{itemize}
\item Exact Accuracy: 52.1\%
\item Adjacent Accuracy: 98.6\%
\item Cohen's $\kappa$: +0.071
\end{itemize}

\textbf{Per-class Recall:} Major 49.8\%, Moderate 55.2\%, Minor 11.4\%.

\textbf{Score-Severity Correlation:} Mean scores by true severity: Major (+0.94), Moderate (+0.09), Minor ($-$0.27). The empirically-derived weights successfully discriminate severity levels despite substantial within-class overlap (std $\approx$ 1.3--1.5).

\newpage
% ============================================================================
% S2: EVIDENCE-BASED CLASSIFIER
% ============================================================================
\section{S2: Evidence-Based Classifier}

\subsection{S2.1 Overview}

The Evidence-Based classifier is a hybrid framework combining the empirically-derived keyword weights from the Rule-Based classifier (Section S1.3) with confidence-based penalization. It uses identical weights derived from DDInter training data but applies a confidence penalty for low-confidence predictions.

\subsection{S2.2 Keyword Scoring}

Uses the same empirically-derived weights as Rule-Based (Table~\ref{tab:empirical_weights}):
\begin{equation}
S = \sum_{k \in \text{keywords}} w_k \cdot \mathbb{1}[k \in \text{description}]
\end{equation}

\subsection{S2.3 Confidence Penalty}

Scores are reduced by 20\% for predictions with base classifier confidence $<$50\%:
\begin{equation}
S_{\text{adjusted}} = 
\begin{cases}
S_{\text{raw}} \times 0.8 & \text{if confidence} < 0.5 \\
S_{\text{raw}} & \text{otherwise}
\end{cases}
\end{equation}

\subsection{S2.4 Percentile-Based Classification}

Final severity is assigned using the same percentile thresholds as Rule-Based, producing our target 4-class distribution:
\begin{equation}
\text{Severity} = 
\begin{cases}
\texttt{Contraindicated} & S \geq P_{96} \text{ (top 4\%)} \\
\texttt{Major} & P_{78} \leq S < P_{96} \text{ (next 18\%)} \\
\texttt{Minor} & S \leq P_{4} \text{ (bottom 4\%)} \\
\texttt{Moderate} & \text{otherwise (middle 74\%)}
\end{cases}
\end{equation}

\subsection{S2.5 Validation Results}

\textbf{Test Set Performance (n=11,150):} Evidence-Based achieves identical performance to Rule-Based (52.1\% exact accuracy, 98.6\% adjacent, $\kappa$=+0.071) because confidence penalization has minimal effect on the empirically-weighted scores.

% ============================================================================
% STATISTICAL ANALYSIS DETAILS
% ============================================================================
\section{S3: Statistical Analysis Details}

\subsection{S3.1 Distribution Divergence Metrics}

\subsubsection{Jensen-Shannon Divergence}
\begin{equation}
D_{JS}(P \| Q) = \frac{1}{2} D_{KL}(P \| M) + \frac{1}{2} D_{KL}(Q \| M)
\end{equation}
where $M = \frac{1}{2}(P + Q)$.

\textbf{Results:}
\begin{itemize}
    \item Original vs.\ Target: $D_{JS} = 0.847$
    \item Recalibrated vs.\ Target: $D_{JS} = 0.000$ (exact match)
    \item Improvement: 100\%
\end{itemize}

\subsubsection{Chi-Square Goodness of Fit}
\begin{equation}
\chi^2 = \sum_{i} \frac{(O_i - E_i)^2}{E_i}
\end{equation}

\textbf{Recalibrated vs.\ Target Distribution:}
\begin{itemize}
    \item $\chi^2 = 12,847$
    \item $df = 3$
    \item $p < 0.001$
\end{itemize}

Note: Statistical significance is expected given the large sample size ($N = 759,774$). The $\chi^2$ test confirms there is still deviation from target, but the practical effect size is small.

\subsection{S3.2 Cohen's Kappa Interpretation}

\begin{table}[h]
\centering
\caption{Cohen's Kappa Interpretation Guide}
\begin{tabular}{ll}
\toprule
\textbf{Kappa Range} & \textbf{Agreement Level} \\
\midrule
$<$ 0.20 & Poor \\
0.21--0.40 & Fair \\
0.41--0.60 & Moderate \\
0.61--0.80 & Substantial \\
0.81--1.00 & Almost perfect \\
\bottomrule
\end{tabular}
\end{table}

Our achieved $\kappa = +0.071$ indicates \textbf{poor agreement} due to class imbalance (DDInter is 73.5% Moderate).

\newpage
% ============================================================================
% DATA AVAILABILITY
% ============================================================================
\section{S4: Data and Code Availability}

\subsection{S4.1 Input Data}
\begin{itemize}
    \item \textbf{Source}: DrugBank v5.1.9 (Creative Commons license)
    \item \textbf{Filter}: Cardiovascular and antithrombotic ATC codes
    \item \textbf{File}: \texttt{data/ddi\_cardio\_or\_antithrombotic\_labeled.csv}
\end{itemize}

\subsection{S4.2 Output Data}
\begin{itemize}
    \item \textbf{Recalibrated predictions}: \texttt{data/ddi\_semantic\_final.csv}
    \item \textbf{Statistics}: \texttt{data/ddi\_semantic\_final.json}
\end{itemize}

\subsection{S4.3 Code}
\begin{itemize}
    \item \textbf{GPU recalibration script}: \texttt{recalibrate\_severity\_gpu.py}
    \item \textbf{Figure generation}: \texttt{publication\_recalibration/generate\_figures\_updated.py}
\end{itemize}

\subsection{S4.4 Reproducibility}
\begin{verbatim}
# Install dependencies
pip install pandas numpy scipy matplotlib torch sentence-transformers

# Run GPU-accelerated semantic recalibration
python recalibrate_severity_gpu.py

# Generate figures
python publication_recalibration/generate_figures_updated.py
\end{verbatim}

\end{document}
